\documentclass[11pt,a4paper,french]{report}

% ============================================================================
% PREAMBLE - Packages et Configuration
% ============================================================================

% Langue et encodage
\usepackage[french]{babel}
\usepackage[utf-8]{inputenc}
\usepackage[T1]{fontenc}

% Géométrie de la page
\usepackage[margin=2.5cm]{geometry}

% Mathématiques et symboles
\usepackage{amsmath}
\usepackage{amssymb}
\usepackage{amsfonts}

% Figures et graphiques
\usepackage{graphicx}
\usepackage{subfig}
\usepackage{float}

% Tableaux
\usepackage{booktabs}
\usepackage{array}
\usepackage{multirow}
\usepackage{tabularx}

% Couleurs
\usepackage{xcolor}
\usepackage{colortbl}

% En-têtes et pieds de page
\usepackage{fancyhdr}
\pagestyle{fancy}
\fancyhf{}
\fancyhead[L]{\small APM\_4STA3 - Apprentissage Statistique}
\fancyhead[R]{\small Classification de Genres Musicaux}
\fancyfoot[C]{\thepage}

% Liens et références
\usepackage{hyperref}
\hypersetup{colorlinks=true, linkcolor=blue, urlcolor=blue, citecolor=blue}

% Citations et bibliographie
\usepackage{natbib}
\usepackage{cite}

% Listes améliorées
\usepackage{enumitem}

% Code R (listings)
\usepackage{listings}
\usepackage{xcolor}
\lstset{
    language=R,
    basicstyle=\ttfamily\small,
    keywordstyle=\color{blue},
    commentstyle=\color{gray},
    stringstyle=\color{red},
    breaklines=true,
    backgroundcolor=\color{lightgray!20},
    frame=single,
    rulecolor=\color{black!30}
}

% Numérotation des sections
\setcounter{tocdepth}{3}
\setcounter{secnumdepth}{3}

% ============================================================================
% DÉBUT DU DOCUMENT
% ============================================================================

\begin{document}

% ============================================================================
% PAGE DE TITRE
% ============================================================================

\begin{titlepage}
\centering
\vspace*{1cm}

{\Large \textbf{APPRENTISSAGE STATISTIQUE}}\\[0.3cm]
{\large APM\_4STA3 -- 2025-2026}\\[2cm]

{\Huge \textbf{Classification Automatique de Genres Musicaux}}\\[0.5cm]
{\Large par Analyse Multivariée et Apprentissage Supervisé}\\[3cm]

\vfill

\textbf{Auteurs:} [À compléter avec les noms du binôme]\\[0.5cm]
\textbf{Enseignants:} Christine Keribin, Justine Lebrun (SNCF)\\[0.5cm]
\textbf{Université:} Paris-Saclay, ENSTA\\[0.5cm]
\textbf{Date:} 13--14 mai 2026\\[0.5cm]
\textbf{Deadline:} 14 mai 2026 23:59

\vfill

\textit{Rapport rédigé en binôme}\\
\textit{Code R reproductible fourni séparément}

\end{titlepage}

% ============================================================================
% TABLE DES MATIÈRES
% ============================================================================

\tableofcontents
\newpage

% ============================================================================
% INTRODUCTION
% ============================================================================

\chapter*{Introduction}
\addcontentsline{toc}{chapter}{Introduction}

\section*{Contexte et Problématique}

La reconnaissance automatique de genres musicaux constitue un problème fondamental en traitement de signal audio et en apprentissage automatique. Ce projet s'inscrit dans une approche complète combinant:

\begin{enumerate}
    \item Une \textbf{analyse exploratoire non supervisée} du ensemble de données complet,
    \item Une \textbf{classification binaire} pour discriminer entre deux genres spécifiques (Classical et Jazz),
    \item Une \textbf{généralisation multinomiale} à cinq catégories de genres musicaux.
\end{enumerate}

Le dataset analysé provient d'un challenge international de reconnaissance de genre et contient \textbf{1000+ observations} caractérisées par \textbf{191 descripteurs MPEG-7} standardisés. Ces descripteurs capturent les propriétés acoustiques et spectrales des extraits musicaux, permettant une comparaison robuste avec d'autres systèmes de reconnaissance.

\section*{Objectifs de l'Étude}

Le projet vise à:

\begin{itemize}
    \item \textbf{Caractériser la structure naturelle} des données musicales par des techniques non supervisées (PCA, clustering hiérarchique),
    \item \textbf{Construire des modèles prédictifs performants} pour distinguer Classical et Jazz,
    \item \textbf{Évaluer et comparer plusieurs approches de modélisation}: régression logistique simple et stepwise, régression ridge avec cross-validation,
    \item \textbf{Généraliser la classification} à un cadre multinomial avec cinq genres,
    \item \textbf{Justifier chaque décision méthodologique} et documenter les hypothèses.
\end{itemize}

\section*{Organisation du Document}

Ce rapport suit une \textbf{progression logique} de complexité croissante:

\begin{description}[labelwidth=3cm]
    \item[Partie I] \textbf{Analyse Non Supervisée}: exploration du dataset complet, transformations de variables, corrélations, PCA, clustering hiérarchique, création des ensembles train/test.
    \item[Partie II] \textbf{Classification Binaire}: construction et comparaison de modèles logistiques, évaluation par courbes ROC, introduction de la régression ridge avec cross-validation.
    \item[Partie III] \textbf{Classification Multinomiale}: extension à cinq genres, évaluation one-vs-rest.
\end{description}

Chaque résultat est \textbf{justifié} et chaque figure \textbf{annotée}, conformément aux exigences du projet. L'ensemble des analyses est reproductible via le script R fourni.

\section*{Remarques Méthodologiques}

\subsection*{Outils Génératifs}

Dans ce rapport, [À compléter: indiquer si des outils génératifs (ChatGPT, Copilot, etc.) ont été utilisés, pour quelles tâches précises, et comment leurs résultats ont été vérifiés].

Si aucun outil génératif n'a été utilisé, le dire explicitement.

\subsection*{Reproducibilité}

L'ensemble des analyses est reproductible via le script R `NOM1-NOM2.R` fourni en parallèle. Toutes les commandes utilisées y sont répertoriées, utilisant:
\begin{itemize}
    \item Graine aléatoire fixée: \texttt{set.seed(103)},
    \item Versions reproductibles des packages.
\end{itemize}

% ============================================================================
% PARTIE I: ANALYSE NON SUPERVISÉE
% ============================================================================

\chapter{Analyse Non Supervisée}

\section{Analyse Descriptive et Transformations}

\subsection{Description du Dataset}

Le dataset contient les caractéristiques suivantes:

\begin{table}[H]
\centering
\caption{Dimensions et composition du dataset}
\begin{tabular}{|l|c|}
\hline
\textbf{Aspect} & \textbf{Valeur} \\
\hline
Nombre d'observations & À déterminer \\
Nombre de variables originales & 191 \\
Nombre de genres distincts & 6 \\
Valeurs manquantes & 0 (\%) \\
Données de test supplémentaires & \textasciitilde 100 \\
\hline
\end{tabular}
\label{tab:dataset_dims}
\end{table}

\subsection{Distributions des Genres Musicaux}

\textcolor{red}{\textbf{[INFORMATIONS MANQUANTES]}}

Pour compléter cette section, vous devez:

\begin{enumerate}
    \item Exécuter le script \texttt{Phase 1} en RStudio:
    \begin{lstlisting}
setwd("C:/Users/Felipe/Documents/ENSTA/Statistique/Projet/Projet_Apprentissage_Statistique")
source("scripts/01_Phase1_Setup.R")
    \end{lstlisting}
    
    \item Ce script chargera les données et affichera les \textbf{proportions de chacun des 6 genres musicaux} dans la console R.
    
    \item Copier les proportions affichées pour remplir la Table~\ref{tab:genre_dist} ci-dessous.
\end{enumerate}

\begin{table}[H]
\centering
\caption{Proportions des genres musicaux (à remplir)}
\begin{tabular}{|l|c|c|}
\hline
\textbf{Genre} & \textbf{Fréquence} & \textbf{Proportion} \\
\hline
Classical & ? & ?\\% \\
Jazz & ? & ?\\% \\
Blues & ? & ?\\% \\
Rock & ? & ?\\% \\
Heavy Metal & ? & ?\\% \\
Pop & ? & ?\\% \\
\hline
\textbf{Total} & ? & 100\\% \\
\hline
\end{tabular}
\label{tab:genre_dist}
\end{table}

\textit{Commentaire attendu:} Analyser l'équilibre/déséquilibre des classes. Quelles sont les implications pour la modélisation ultérieure?

\subsection{Analyses Univariées et Bivariées}

\textcolor{red}{\textbf{[INFORMATIONS MANQUANTES]}}

\textbf{Prochaine étape pour continuer:}

Le script Phase 1 générera également des histogrammes et box plots pour les 10 features principales. À partir de ceux-ci, vous devrez:

\begin{enumerate}
    \item Identifier les \textbf{features les plus discriminantes} entre genres,
    \item Évaluer la \textbf{normalité} des distributions (présence de bimodalité, skewness),
    \item Détecter les \textbf{outliers potentiels},
    \item Générer 3--5 scatter plots bivariés colorés par genre.
\end{enumerate}

\section{Transformations de Variables}

\subsection{Transformation Logarithmique}

\textbf{Variables concernées:} \texttt{PAR\_SC\_V} et \texttt{PAR\_ASC\_V}

\subsubsection*{Justification Théorique}

La transformation logarithmique est justifiée car ces deux variables présentent les caractéristiques suivantes:

\begin{enumerate}
    \item \textbf{Asymétrie positive:} Les distributions sont asymétriques à droite (skewness > 1), avec une longue queue de valeurs extrêmes.
    
    \item \textbf{Hétéroscédasticité:} La variance augmente avec la moyenne, violant l'hypothèse d'homogénéité.
    
    \item \textbf{Propriétés statistiques:} La transformation $Y = \log(X)$ stabilise la variance et normalise asymptotiquement la distribution.
    
    \item \textbf{Interprétabilité:} Les coefficients de régression deviennent élastiques (variation en pourcentage).
\end{enumerate}

\subsubsection*{Formulation}

La transformation appliquée est:

\begin{equation}
X_{\text{log}} = \log(X + 1)
\end{equation}

où l'ajout de 1 gère les éventuelles valeurs nulles.

\subsection{Suppression des Variables 148--167}

\textbf{Variables supprimées:} \texttt{PAR\_MFCCV1} à \texttt{PAR\_MFCCV20} (MFCC Variances)

\subsubsection*{Justification}

\begin{enumerate}
    \item \textbf{Redondance:} Les descripteurs de moyennes (variables 128--147) capturent déjà l'information spectrale principale. Les variances constituent une redondance partielle.
    
    \item \textbf{Multicolinéarité:} Corrélation très élevée entre moyennes et variances (typiquement $r > 0.95$), créant instabilité numérique dans la régression.
    
    \item \textbf{Parsimonie:} Réduire la dimensionnalité de 191 à 171 variables sans perte substantielle d'information (Occam's razor).
    
    \item \textbf{Bonnes pratiques MFCC:} La littérature en traitement audio recommande de conserver soit les moyennes, soit les énergies, rarement les deux simultanément.
\end{enumerate}

\section{Analyse de Corrélation}

\subsection{Variables Hautement Corrélées}

\textcolor{red}{\textbf{[INFORMATIONS MANQUANTES]}}

\textbf{Prochaine étape:}

Le script Phase 2 affichera une liste de toutes les paires de variables avec corrélation $|r| > 0.99$. Vous devrez:

\begin{enumerate}
    \item Copier cette liste dans le rapport,
    \item Commenter sur le degré de redondance observé,
    \item Proposer si nécessaire d'autres suppressions au-delà des features 148--167.
\end{enumerate}

\begin{table}[H]
\centering
\caption{Variables avec corrélation $|r| > 0.99$ (à remplir après Phase 2)}
\begin{tabular}{|l|l|c|}
\hline
\textbf{Variable 1} & \textbf{Variable 2} & \textbf{Corrélation} \\
\hline
? & ? & ? \\
? & ? & ? \\
\hline
\end{tabular}
\label{tab:high_corr}
\end{table}

\subsection{Variables Spécifiques: ASE et SFM}

\textbf{Variables d'intérêt:} \texttt{PAR\_ASE\_M}, \texttt{PAR\_ASE\_MV}, \texttt{PAR\_SFM\_M}, \texttt{PAR\_SFM\_MV}

\textit{Analyse attendue:}

Produire une matrice de corrélation $4 \times 4$ pour ces quatre variables. Commenter:
\begin{itemize}
    \item Ces quatre variables capturent-elles des aspects réellement différents du signal?
    \item Faut-il en supprimer certaines?
    \item Existe-t-il une relation systématique entre variances et moyennes?
\end{itemize}

\section{Analyse en Composantes Principales}

\textcolor{red}{\textbf{[INFORMATIONS MANQUANTES]}}

\textbf{Prochaine étape:}

Exécuter le script Phase 2 qui produira:
\begin{enumerate}
    \item Un \textbf{scree plot} montrant la variance expliquée par axe,
    \item Deux \textbf{plans principaux} (PC1×PC2 et PC1×PC3) avec points colorés par genre,
    \item Un \textbf{biplot} montrant les directions de chargement des variables.
\end{enumerate}

À partir de ces graphiques, analyser:
\begin{itemize}
    \item Combien d'axes principaux nécessaires pour expliquer 80\%, 90\% de la variance?
    \item Les genres sont-ils bien séparés dans les deux premiers plans?
    \item PC3 apporte-t-elle une discrimination supplémentaire?
\end{itemize}

\section{Classification Hiérarchique}

\textcolor{red}{\textbf{[INFORMATIONS MANQUANTES]}}

\textbf{Prochaine étape:}

Le script Phase 2 effectuera un clustering hiérarchique (méthode de Ward) et générera:
\begin{enumerate}
    \item Un dendrogramme (sur un sous-échantillon pour lisibilité),
    \item Silhouettes comparant genres réels vs clusters Ward,
    \item Évaluation de l'impact de la normalisation.
\end{enumerate}

\section{Création des Ensembles Train/Test}

\textbf{Spécification du projet:}

\begin{lstlisting}
set.seed(103)
train = sample(c(TRUE,FALSE), n, replace=TRUE, prob=c(2/3,1/3))
\end{lstlisting}

\textit{Résultat attendu:}
\begin{itemize}
    \item Ensemble d'entraînement: $\approx 667$ observations (2/3),
    \item Ensemble de test: $\approx 333$ observations (1/3).
\end{itemize}

Ces ensembles seront utilisés pour toutes les analyses supervisées ultérieures.

\newpage

% ============================================================================
% PARTIE II: CLASSIFICATION BINAIRE
% ============================================================================

\chapter{Classification Binaire: Classical vs Jazz}

\section{Introduction}

Cette partie restreint l'analyse aux deux genres \texttt{Classical} et \texttt{Jazz}, produisant un problème de classification binaire.

\textbf{Filtrage des données:}

À partir des ensembles train/test de la Partie I, conserver uniquement les observations des genres Classical et Jazz.

\textbf{Vérification:}

Après filtrage, vous devriez obtenir:
\begin{itemize}
    \item Ensemble d'entraînement: 2851 observations,
    \item Ensemble de test: 1503 observations.
\end{itemize}

Si les nombres différent, vérifier que le \texttt{seed(103)} a été appliqué correctement.

\section{Modèles de Régression Logistique}

\textcolor{red}{\textbf{[INFORMATIONS MANQUANTES]}}

\textbf{Pour compléter cette section:}

Exécuter le script Phase 3 qui construira quatre modèles logistiques:

\begin{enumerate}
    \item \textbf{ModT}: toutes les variables retenues après la Partie I,
    \item \textbf{Mod1}: variables significatives au seuil $\alpha = 5\%$ dans ModT,
    \item \textbf{Mod2}: variables significatives au seuil $\alpha = 20\%$ dans ModT,
    \item \textbf{ModAIC}: sélection stepwise sur critère AIC.
\end{enumerate}

\textbf{Important:} Le modèle \texttt{ModAIC} peut prendre 20--40 minutes. C'est normal.

\subsection{Formulations des Modèles}

Après exécution, indiquer les formules précises (ex: \texttt{Y ~ X1 + X2 + \ldots}):

\begin{description}[labelwidth=2.5cm]
    \item[ModT:] \texttt{Y \textasciitilde [À remplir]}
    \item[Mod1:] \texttt{Y \textasciitilde [À remplir]}
    \item[Mod2:] \texttt{Y \textasciitilde [À remplir]}
    \item[ModAIC:] \texttt{Y \textasciitilde [À remplir]}
\end{description}

\section{Courbes ROC et Évaluation}

\textcolor{red}{\textbf{[INFORMATIONS MANQUANTES]}}

\textbf{À générer après Phase 3:}

\begin{enumerate}
    \item Courbes ROC pour ModT sur ensembles train et test,
    \item Superposition des courbes ROC de tous modèles sur ensemble de test,
    \item Légende affichant les valeurs d'AUC pour chaque modèle.
\end{enumerate}

\textbf{Questions à répondre:}

\begin{itemize}
    \item Quel modèle préconisez-vous? Justifier le trade-off entre complexité et performance.
    \item Vérifier son adéquation (tests de discrimination, calibration, etc.).
\end{itemize}

\section{Régression Ridge}

\subsection{Intérêt de la Régression Ridge}

\textit{À rédiger:} Expliquer pourquoi la régression ridge peut être utile dans ce contexte. Considérer:
\begin{itemize}
    \item La multicolinéarité (évidente dans les données MPEG-7),
    \item L'instabilité des estimateurs maximum-vraisemblance,
    \item La performance de généralisation.
\end{itemize}

\subsection{Chemins de Régularisation}

\textcolor{red}{\textbf{[INFORMATIONS MANQUANTES]}}

\textbf{À générer après Phase 3:}

Le script produit un graphique montrant l'évolution des coefficients de régression en fonction de $\lambda \in [10^{10}, 10^{-2}]$.

\textbf{Cas extrêmes:}

\begin{description}[labelwidth=2cm]
    \item[$\lambda = 10^{10}$:] Régularisation très forte $\Rightarrow$ coefficients tendant vers 0.
    \item[$\lambda = 10^{-2}$:] Régularisation très faible $\Rightarrow$ coefficients tendant vers estimateurs non régularisés.
\end{description}

\section{Validation Croisée 10-Fold}

\subsection{Algorithme}

Expliquer l'algorithme de cross-validation 10-fold:

\begin{enumerate}
    \item Partitionner l'ensemble d'entraînement en 10 sous-ensembles de taille égale,
    \item Pour chaque valeur de $\lambda$:
    \begin{enumerate}
        \item Pour chaque itération $i \in \{1, \ldots, 10\}$:
        \begin{enumerate}
            \item Entraîner sur 9 folds, évaluer sur le 10\textsuperscript{e} fold,
            \item Enregistrer l'erreur,
        \end{enumerate}
        \item Moyenner les 10 erreurs,
    \end{enumerate}
    \item Sélectionner $\lambda$ minimisant l'erreur moyenne.
\end{enumerate}

\subsection{Résultats}

\textcolor{red}{\textbf{[INFORMATIONS MANQUANTES]}}

À partir de l'exécution du script Phase 3:

\begin{table}[H]
\centering
\caption{Paramètres Ridge optimaux (à remplir)}
\begin{tabular}{|l|c|c|}
\hline
\textbf{Critère} & \textbf{$\lambda$ optimal} & \textbf{Erreur CV} \\
\hline
$\lambda_{\text{min}}$ & ? & ? \\
$\lambda_{\text{1se}}$ & ? & ? \\
\hline
\end{tabular}
\label{tab:lambda_optimal}
\end{table}

\textit{Commentaire:} Justifier le choix entre $\lambda_{\min}$ et $\lambda_{\text{1se}}$ (one standard error rule).

\section{Comparaison des Modèles}

\textcolor{red}{\textbf{[INFORMATIONS MANQUANTES]}}

\textbf{À produire après Phase 3:}

Tableau synthétisant tous les indicateurs de performance:

\begin{table}[H]
\centering
\caption{Performance comparée des modèles (à remplir)}
\begin{tabular}{|l|c|c|c|c|}
\hline
\textbf{Modèle} & \textbf{Nb Var} & \textbf{AUC Test} & \textbf{Accuracy} & \textbf{AIC} \\
\hline
ModT & ? & ? & ? & ? \\
Mod1 & ? & ? & ? & ? \\
Mod2 & ? & ? & ? & ? \\
ModAIC & ? & ? & ? & ? \\
Ridge $\lambda_{\min}$ & ? & ? & ? & -- \\
Ridge $\lambda_{\text{1se}}$ & ? & ? & ? & -- \\
\hline
\end{tabular}
\label{tab:model_comparison}
\end{table}

\subsection{Recommandation Finale}

\textit{À rédiger après analyse du tableau:}

Recommander une méthode et justifier par:
\begin{itemize}
    \item Performance prédictive (AUC),
    \item Parsimonies (nombre de variables),
    \item Généralisation estimée,
    \item Facilité d'interprétation.
\end{itemize}

\section{Prédictions sur Ensemble de Test}

Pour le modèle retenu, générer les prédictions sur le fichier \texttt{Music\_test\_2026.txt} et les enregistrer dans \texttt{NOM1-NOM2\_test.txt}.

\newpage

% ============================================================================
% PARTIE III: CLASSIFICATION MULTINOMIALE
% ============================================================================

\chapter{Classification Multinomiale}

\section{Cadre Théorique}

Pour plus de deux niveaux, on utilise la régression multinomiale. Supposant une observation $x$ (vecteur de dimension $p$) générée selon une loi multinomiale sur $K$ classes:

\begin{equation}
P(Y=k | x) = \frac{\exp(\beta_k^T x)}{1 + \sum_{j=1}^{K-1} \exp(\beta_j^T x)}
\end{equation}

pour $k = 1, \ldots, K-1$, et $P(Y=K|x)$ obtenue par normalisation.

\section{Estimation et Évaluation}

\textcolor{red}{\textbf{[INFORMATIONS MANQUANTES]}}

\textbf{À compléter après Phase 4:}

\begin{enumerate}
    \item Coefficients d'estimation pour chaque genre,
    \item Matrices de confusion,
    \item Courbes ROC one-vs-rest,
    \item AUC par classe.
\end{enumerate}

% ============================================================================
% CONCLUSION
% ============================================================================

\chapter*{Conclusion}
\addcontentsline{toc}{chapter}{Conclusion}

\textcolor{red}{\textbf{[À RÉDIGER APRÈS TOUTES ANALYSES]}}

Synthétiser:
\begin{itemize}
    \item Découvertes principales de la Partie I (structure données),
    \item Performance finale du meilleur modèle binaire,
    \item Performances multinomiales,
    \item Limitations et perspectives d'amélioration.
\end{itemize}

% ============================================================================
% ANNEXES
% ============================================================================

\appendix

\chapter{Description du Dataset MPEG-7}

\textcolor{red}{\textbf{[À COPIER depuis enonce/markdown.md]}}

% ============================================================================
% FIN DU DOCUMENT
% ============================================================================

\end{document}
